\section{System Architecture and Analytical Methodology}
\label{sec:system}

Figure~\ref{fig:architecture} shows the workflow as implemented. An analyst of the
monitored organization uploads one file per evidence category through an HTTP API;
validation stores per-record provenance; a separate worker process claims queued
runs under a lease and executes the analysis; a supervisor reviews and decides; and
the decision is finalized and bound.

\begin{figure}[!htbp]
\centering
% Figure 1 -- SAT-SA workflow (native TikZ; edit labels directly).
\begin{tikzpicture}[node distance=3.5mm and 6mm, every node/.style={font=\scriptsize},
  box/.style={draw, rounded corners=1pt, align=center, minimum width=40mm, minimum height=6mm},
  arr/.style={-{Stealth[length=1.8mm]}}]
\node[box] (sub) {Periodic evidence submission\\(alerts, cases, steps, escalations,\\dispositions, assets)};
\node[box, below=of sub] (val) {Validation and normalization\\(per-record provenance, all-or-nothing)};
\node[box, below=of val, minimum height=13mm] (ana) {Analytical workers (\V{CODE-workers}, rule-based)\\execution $\cdot$ negative space\\deviation $\cdot$ peer $\cdot$ coverage $\cdot$ context};
\node[box, below=of ana] (risk) {Fusion into \V{CODE-dimensions} risk dimensions\\and prioritization};
\node[box, right=14mm of sub] (find) {Evidence-cited findings\\and recommendations};
\node[box, below=of find] (rev) {Human supervisory review\\(review interrupt)};
\node[box, below=of rev] (dec) {Authorized decision};
\node[box, below=of dec] (trust) {TRUST-SAT finalization: SHA3-256\\record, \V{CODE-signature} receipt, ledger};
\node[box, below=of trust] (ver) {Verification and audit};
\draw[arr] (sub) -- (val); \draw[arr] (val) -- (ana); \draw[arr] (ana) -- (risk);
\draw[arr] (risk.east) -| ++(5mm,0) |- (find.west);
\draw[arr] (find) -- (rev); \draw[arr] (rev) -- (dec); \draw[arr] (dec) -- (trust); \draw[arr] (trust) -- (ver);
\end{tikzpicture}

\caption{SAT-SA supervisory workflow along the evaluated path. The analytical-worker
stage is the evaluated subset of the registered agents (Table~\ref{tab:architecture}).
All analytics are deterministic and rule-based; orchestration (direct worker or
LangGraph) sequences the stages and pauses for the human decision.}
\label{fig:architecture}
\end{figure}

\subsection{Architecture layers}
The implementation registers \V{CODE-registry} components that it calls
\emph{agents} (Table~\ref{tab:architecture}): \V{CODE-registry-satsa} supervisory agents
and \V{CODE-registry-mlops} MLOps agents retained from the platform SAT-SA was built on.
An agent here is a registered specification (purpose, inputs, outputs, evidence types)
that maps to deterministic code; none is a language-model or autonomous agent. Eleven
supervisory agents cover the analysis families and group the \V{CODE-workers} default
analytical workers evaluated in this paper; the other twelve cover the pipeline around
them, from ingestion and normalization through correlation, risk scoring,
prioritization and recommendation to review, trust and provenance, reporting and
validation. The reported experiments exercise the analytical workers and the
pipeline from ingestion to trust; evidence assembly, meta-audit, report generation,
the calibration-oriented validation agent and the retained MLOps agents are not
evaluated. The count of \V{CODE-workers} therefore refers to the analytical workers
only, not to the whole architecture.

\begin{table}[!htbp]
\caption{SAT-SA architecture inventory as implemented. ``Agent'' is the
implementation's name for a registered component; each maps to deterministic code, and
none is a language-model or autonomous agent. Only the rows marked ``yes'' are exercised
by the reported experiments.}
\label{tab:architecture}
\centering
\small
\setlength{\tabcolsep}{3pt}
\begin{tabular}{@{}L{3.5cm}L{1.0cm}L{5.1cm}L{1.8cm}@{}}
\toprule
Layer & Count & Members & Evaluated \\
\midrule
Registered agents (agent registry) & \V{CODE-registry} & \V{CODE-registry-satsa} supervisory agents and \V{CODE-registry-mlops} retained MLOps agents & in part \\
\quad Supervisory analytical agents & 11 & one agent per analysis family; they group the \V{CODE-workers} default analytical workers (the execution-gap agent groups six detectors) & yes \\
\quad Supervisory pipeline agents & 12 & ingestion, normalization, correlation, risk scoring, prioritization, recommendation, review workflow, trust--provenance, evidence assembly, meta-audit, report generation, validation & ingestion to trust: yes; last four: no \\
\quad Retained MLOps agents & \V{CODE-registry-mlops} & data, performance, security, quantum security, red team, training optimization, incident response, governance, optimization (inherited platform) & no \\
Default analytical workers & \V{CODE-workers} & listed in Tables~\ref{tab:workers} and~\ref{tab:workers-b} & yes \\
Orchestration graph nodes & \V{CODE-graph-nodes} & readiness, analysis, recommendations, human review, trust boundary (LangGraph mode) & yes \\
\bottomrule
\end{tabular}
\end{table}


\subsection{Analytical workers}
The \V{CODE-workers} default analytical workers fall into five functional groups.
\emph{Execution and process rules} (fast closure, acknowledgement without
investigation, critical alert without escalation, repeated shallow investigation,
recurrence without remediation, potential metric gaming) check individual records and
short record chains against expected handling. \emph{Negative-space analysis} detects
the absence of expected work: cases without investigation steps, missing escalations
or dispositions, critical assets without alerts; these are rule-encoded conformance
checks. \emph{Deviation analysis} comprises within-entity anomalies, drift and the peer
rule of Eq.~\eqref{eq:peer}. \emph{Coverage and completeness} (coverage gap, evidence
completeness) flag unmonitored assets and missing evidence categories, and
entity--asset resolution links records to the assets they concern. \emph{Context workers} (workflow reconstruction, case similarity,
cross-entity insights) reconstruct sequences and clusters; their signals also feed risk
dimensions. Tables~\ref{tab:workers} and~\ref{tab:workers-b} list all \V{CODE-workers} workers with the
dimension each feeds and its confidence rule. Because a rule's dimension is assigned by
name prefix, the escalation-discipline and investigation-quality dimensions are fed only
by workflow reconstruction (and by keyword-matching drift or cross-entity rules), while
critical alerts without escalation count towards the execution gap. No worker uses a
language model or a learned model.

\begin{table}[!htbp]
\caption{Default analytical workers evaluated in the reported experiments (part 1 of
2; the analytical subset of the registered agents, Table~\ref{tab:architecture}): risk
dimension fed and confidence components, $c_f=\min(a_f,e_f[,p_f])$
(Eq.~\eqref{eq:confidence}). $q$ = qualifying records; $N$ = records of the relevant
category; $\delta$ = effect size clipped to $[0,1]$; $z$ = deviation in MADs. All
constants are fixed in code. Dimensions: EG execution gap, NS negative space, AN
anomaly, PD peer deviation, DG detection gap, IQ investigation quality, ED escalation
discipline; kw = assigned by rule-name keyword (default AN).}
\label{tab:workers}
\centering
\small
\setlength{\tabcolsep}{3pt}
\begin{tabular}{@{}L{2.7cm}L{0.9cm}L{4.4cm}L{3.5cm}@{}}
\toprule
Worker & Dim. & Analytical support $a_f$ & Evidence completeness $e_f$ (peer $p_f$) \\
\midrule
Fast closure & EG & $0.4+0.6\,\delta$, $\delta=1-\text{median}/\text{limit}$ & $q/N$ (alerts of that severity) \\
Ack.\ without investigation & EG & $0.5+0.4\,q/N$ & $q/N$ (alerts) \\
Critical without escalation & EG & $0.8$ & $q/N$ (alerts) \\
Repeated investigation & EG & $0.4+0.5\,\delta$ & $\min(1,N_{\text{steps}}/20)$ \\
Recurrence without remediation & EG & $0.7$ & $q/N$ (cases) \\
Potential metric gaming & EG & $0.4+0.4\,\delta$ & $\min(1,N/20)$ \\
Negative space & NS & $0.7$ steps; $0.6$ escalation, monitoring; $0.5$ disposition; $0.4$ low activity & $1$ if category submitted, else $0$ \\
Anomaly & AN & $0.4+0.5\min(1,z/3)$ & $1$ \\
Peer benchmark & PD & $0.4+0.1\min(1,|z|/3)$ & $1$; $p_f=\min(1,n_{\text{peers}}/10)$ \\
Drift & kw & $0.4+0.6\min(1,|\Delta_{\text{rel}}|)$ & $0.9/0.6/0.3$ (${\ge}3$/${\ge}1$/$0$ refs) \\
\bottomrule
\end{tabular}
\end{table}

\begin{table}[!htbp]
\caption{Default analytical workers evaluated in the reported experiments (part 2 of
2); abbreviations as in Table~\ref{tab:workers}.}
\label{tab:workers-b}
\centering
\small
\setlength{\tabcolsep}{3pt}
\begin{tabular}{@{}L{2.7cm}L{0.9cm}L{4.4cm}L{3.5cm}@{}}
\toprule
Worker & Dim. & Analytical support $a_f$ & Evidence completeness $e_f$ (peer $p_f$) \\
\midrule
Coverage gap & DG & $0.4+0.6\,g$ ($g$ unmonitored critical-asset share) & $0.9$ if alerts link assets, else $0.4$ \\
Cross-entity insights & kw & $0.4+0.6\min(1,\pi)$ ($\pi$ prevalence) & $0.9$ if ${\ge}3$ entities, else $0.6$; $p_f=\min(1,\pi)$ \\
Case similarity & AN & $0.5+0.5\,s$ ($s$ cluster share) & $0.7$ \\
Evidence completeness & AN & $0.4+0.6\,m$ ($m$ missing share) & $1-m$ \\
Workflow reconstruction & ED, IQ, AN & $0.85$ / $0.8$ / $0.6$ by rule & $1$ \\
Entity--asset resolution & AN & $0.6$ & $1$ \\
\bottomrule
\end{tabular}
\end{table}


\subsection{Fusion, prioritization and recommendations}
Signal findings are fused by Eq.~\eqref{eq:risk} into an entity risk decomposed by
dimension, so that a supervisor sees which findings produce which part of the score.
Entities are ordered by Eq.~\eqref{eq:priority}, and each finding family maps to a
recommended supervisory action.

\subsection{Human supervisory review}
Runs can require review. The system never decides: only a supervisor, or an
organization administrator, can record the decision, after which the worker resumes.
Automation assists; the supervisor retains terminal authority.

\subsection{Stateful orchestration}
The same stages run in two modes. In \emph{direct} mode the worker pauses the run after
analysis and recommendations. In \emph{graph} mode a LangGraph~\cite{langgraph2026docs}
graph with \V{CODE-graph-nodes} nodes (readiness, analysis, recommendations, human
review, trust boundary) checkpoints its state and interrupts at the review node.
LangGraph provides checkpointing, interrupt/resume and recovery; it does not replace or
alter the deterministic workers, and both modes produce identical findings and risk
(Section~\ref{sec:results}).

\subsection{Integrity layer}
After the decision, finalization binds the reviewed state as defined in
Section~\ref{sec:formulation}: canonical serialization, a SHA3-256
commitment~\cite{nist2015fips202}, an \V{CODE-signature}~\cite{nist2024fips204}
signature, a stored receipt and a hash-chained ledger entry in the spirit of
tamper-evident logging~\cite{schneier1999logs,crosby2009logging}. The threat model is
after-the-fact alteration of the supervisory record by someone with write access to
the database or ledger file. Out of scope are compromise of the signing key, of the
running process before finalization or of the operating system, and false records
submitted by the monitored organization: TRUST-SAT binds what was analysed and
decided, not whether the submitted records are true. An attacker who can rewrite the
database and re-sign with the service key, or substitute a key pair consistently, is
not detected unless the public key is pinned out of band.
